\section{Conclusion}
\label{sec:conclusion}
Applied to the regression problem, the neural network
showed a slightly lower test mse than the linear models, which is an
impressive result. Although this is impressive, and shows 
that the neural network works, the data cost of the neural network,
pushes the weight in favor of the linear models for solving
this particular problem, the runge function on normal randomly distributed data.
When applying the same model structure to 
the classification problem, with the only significant difference being
the output activation function, the cost function
and the backpropagation. The model showed a great accuracy
in predicting unknown digits, when trained on only 1000 digits.
With code optimization or a computer with more power, the 
model could be trained on more data, and thus achieve a higher overall accuracy.
A large portion of this project was 
used in interpreting the results of the
GridSearchCV class created by
Scikit learn, using different 
cross validation methods, and
a better understanding of this class,
would have further improved this project
and the results of the training and
cross validation. 
A large part of the success in this process
can be attributed to stochastic gradient descent,
the backpropagation alghoritm and the 
gradient updating method adam. Great methods
for randomizing and increasing the training data, 
calculating gradients with a low data cost and for 
optimizing the gradients in an effective manner.